\documentclass{article}
\usepackage[T1]{fontenc}
\usepackage{lmodern}
\usepackage{graphicx}
\usepackage{amsmath,amssymb}
\usepackage[a4paper,margin=2cm]{geometry}
\usepackage[absolute,overlay]{textpos}
\setlength{\TPHorizModule}{1mm}
\setlength{\TPVertModule}{1mm}
\usepackage[indonesian]{babel}
\usepackage{hyperref}
\setlength{\emergencystretch}{2em}
% unit: Halaman 58

\begin{document}

% === Halaman 58 (python/native) ===

(ii) a  b = (x6 + x4 + x2 + x + 1 )  (x7 + x + 1 )  

                 = x13 + x7 + x6 + x11 + x5 + x4 + x9 + x3 + x2 + x8 + x2 + x + x7 + x + 1 (mod 2)


    = x13 + x11 + x9 + x8 + 2x7 + x6 + x5 + x4 + x3 + 2x2 + 2x + 1 (mod 2) 


    = x13 + x11 + x9 + x8 + 0x7 + x6 + x5 + x4 + x3 + 0x2 + 0x + 1 

                = x13 + x11 + x9 + x8 + x6 + x5 + x4 + x3 + 1  


   Karena m = 8, maka hasil perkalian di atas direduksi menjadi derajat < 8 oleh 

 irreducible polynomial . Algoritma Rijndael menggunakan irreducible polynomial

   f(x) = x8 +  x4 +  x3 + x + 1 

  x13 + x11 + x9 + x8 + x6 + x5 + x4 + x3 + 1 (mod f(x)) = x7 + x6 + 1  = 11000001




  a  b = 11000001 = ‘C1’

\end{document}
